% TOP_PROBLEM: {"id": 30006961, "title": "Persistence Conjecture for Weakly Reversible Mass-Action Systems", "category_id": 11, "category": {"id": 11, "name": "dynamical_systems", "display_name": "Dynamical Systems", "description": "Problems about long-term behavior of deterministic systems, Hamiltonian dynamics, and periodic orbits.", "slug": "dynamical-systems", "order_index": 11, "created_at": "2026-07-31T15:26:25.671Z"}, "status": "open"}
% FIRST_POSTED: 2026-09-27
% !TeX program = lualatex
% ATTEMPT_STATUS: unresolved
% Model: GPT-6 Astra Ultra; continuation dated 27 September 2026.
\documentclass[11pt]{article}
\usepackage[margin=25mm]{geometry}
\usepackage{fontspec}
\setmainfont{Latin Modern Roman}
\usepackage{amsmath,amssymb,mathtools}
\usepackage{unicode-math}
\setmathfont{Latin Modern Math}
\usepackage{xurl,hyperref}
\hypersetup{colorlinks=true,urlcolor=blue}
\setcounter{secnumdepth}{0}
\setlength{\parindent}{0pt}
\setlength{\parskip}{5pt}
\setlength{\emergencystretch}{3em}
\begin{document}
\section*{\texorpdfstring{Persistence Conjecture for Weakly Reversible Mass-Action Systems}{Persistence Conjecture for Weakly Reversible Mass-Action Systems}}\label{30006961}
\subsection{Definitions and mathematical statement}
A reaction network has reactions $y\to y'$ between vectors $y,y'\in{\mathbb{Z}}_{\ge0}^d$, with positive rate constants $k_{y\to y'}$. Its mass-action equation is
\[
 \dot x=\sum_{y\to y'}k_{y\to y'}x^y(y'-y),
 \qquad x^y=\prod_{i=1}^dx_i^{y_i}.
\]
It is weakly reversible if every directed reaction belongs to a directed cycle of complexes. Let $S$ be the real span of all reaction vectors $y'-y$. The stoichiometric compatibility class through $x_0>0$ is $(x_0+S)\cap{\mathbb{R}}_{>0}^d$.

Assuming this class is bounded, the conjecture asks that every trajectory starting at $x_0$ be persistent:
\[
 \liminf_{t\to\infty}x_i(t)>0\qquad(i=1,\ldots,d).
\]
Equivalently in the bounded setting, the omega-limit set avoids the boundary of the nonnegative orthant. Mere positivity at each finite time is weaker; it allows asymptotic approach to extinction.
\par\smallskip\textit{Formulation and concept references:} \href{https://reaction-networks.net/wiki/Persistence_and_the_Global_Attractor_Conjecture}{[S1]}.

\subsection{Short English statement}
In a weakly reversible mass-action network with bounded compatibility class, must every initially present species remain bounded away from extinction at late times?
\subsection{Research attempt}

\paragraph{Scope and a significant recent source update (27 September 2026).}
This GPT-6 Astra Ultra attempt studies the bounded-class formulation
printed above. The rate constants are arbitrary positive constants and
weak reversibility is a condition on the directed graph of complexes.
It does not assert complex balance for the chosen rates. Persistence is
a lower bound along each individual positive trajectory, not a single
bound uniform over all initial conditions in the class.

There is a recent claim directly relevant to the selected scope. Craciun's
\emph{Toric Differential Inclusions and a Proof of the Global Attractor
Conjecture}, arXiv:1501.02860v3, was revised on 23 September 2026 [E1].
The final remark in Section 8 explicitly claims persistence of every
bounded trajectory of a variable-$k$ weakly reversible system. Constant
positive rates are a special case, and a bounded compatibility class
forces its trajectory to be bounded. Thus the stated claim covers this
entry; it is not merely the complex-balanced global-attractor assertion.
The manuscript's Theorem A / Theorem 3.1 supplies the embedding into toric
differential inclusions, and its higher-dimensional invariant-region
construction supplies the asserted boundary exclusion. The version note
describes a major revision of the general-dimensional extension.

The source statement and version date were verified, but the full
high-dimensional proof has not been independently audited in this
attempt. Accordingly, the retained catalogue statement is not accompanied
by an unqualified assertion that its frontier is unchanged or that the
recent claim has been independently certified. The status ``unresolved''
below describes the independent argument developed here. Older primary
work proves a siphon/conservation criterion [E2] and persistence for
bounded trajectories with two-dimensional stoichiometric subspace [E3].
The proofs below reconstruct elementary parts of that established theory.

The attack is to reduce boundary accumulation to special coordinate
faces, exclude faces carrying suitable conservation laws, and then examine
the faces not excluded that way. In one stoichiometric dimension a
monomial-ordering argument finishes the proof. An explicit reversible
example shows why that ordering does not reduce to a conservation-law
argument alone, while another example separates weak reversibility from
complex balance.

\paragraph{Bounded classes, positive conservation, and finite-time existence.}
Let $\mathcal C=(x_0+S)\cap\mathbb R_{>0}^d$, where $x_0>0$ and
$\mathcal C$ is bounded. Its closure is
$(x_0+S)\cap\mathbb R_{\geq0}^d$: every point in the latter is a limit
of its convex combinations with $x_0$. A vector $w\in S^\perp$ gives a
conservation law $w\cdot x(t)=w\cdot x_0$.

\textbf{Lemma 481.1.} Boundedness of $\mathcal C$ is equivalent to
$S\cap\mathbb R_{\geq0}^d=\{0\}$, and implies that there exists
$w\in S^\perp$ with every coordinate strictly positive.

\emph{Proof.}
A nonzero $v\in S\cap\mathbb R_{\geq0}^d$ gives the unbounded ray
$x_0+t v\in\mathcal C$. Conversely, if $x_n\in\mathcal C$ is
unbounded, pass to a subsequence of
$(x_n-x_0)/\|x_n-x_0\|$ converging to a unit vector $v$.
It belongs to $S$ and has nonnegative coordinates, a contradiction.

For the strictly positive conservation vector, consider the simplex
$D=\{z\geq0:\sum_i z_i=1\}$ and its orthogonal projection $K$ onto
$S^\perp$. It is compact and convex and does not contain zero; otherwise
some element of $D$ would lie in $S$. Let $w\in K$ have minimum
Euclidean norm. Convexity and the minimum property imply
$w\cdot(z-w)\geq0$ for every $z\in K$, by differentiating the squared
norm along the segment from $w$ to $z$. Apply this to the projection of
each coordinate vector $e_i$. Since $w\in S^\perp$, one obtains
$w_i\geq\|w\|^2>0$ for every $i$.\hfill$\square$

The mass-action vector field is polynomial and hence locally Lipschitz.
At a point with $x_i=0$, every reaction consuming species $i$ has rate
zero because its reactant exponent $y_i$ is at least one. Contributions
from reactions not consuming $i$ are nonnegative in that coordinate.
The nonnegative orthant is therefore forward invariant.

The positive conservation vector bounds every coordinate by
$M_i=(w\cdot x_0)/w_i$. On the resulting compact box, each negative
term in $\dot x_i$ contains a factor $x_i$, and all remaining factors
are bounded. There is a finite constant $C_i$ with
\[
 \dot x_i\geq-C_i x_i,
 \qquad x_i(t)\geq x_i(0)e^{-C_i t}>0.
\]
No solution can blow up in finite time because it stays in the compact
box; the local existence theorem therefore extends it to all $t\geq0$.
These facts justify taking a nonempty compact omega-limit set. The
displayed lower bound tends to zero as $t\to\infty$, so it is not a
proof of persistence.

\paragraph{Boundary accumulation forces a siphon.}
A nonempty set $Z\subseteq\{1,\ldots,d\}$ is a \emph{siphon} if
every reaction whose product contains a species in $Z$ has a reactant
containing at least one species in $Z$. Equivalently,
\[
 (\exists i\in Z:y'_i>0)\quad\Longrightarrow\quad
                         (\exists j\in Z:y_j>0)
 \qquad\text{for every reaction }y\to y'.
\]
This is also called a semilocking set. The relative coordinate face
associated to $Z$ consists of points with $x_i=0$ exactly for $i\in Z$.
If $Z$ is a siphon, the closed face with all those coordinates zero is
forward invariant: every reaction producing one of them has zero rate
there.

\textbf{Lemma 481.2.} If a bounded positive trajectory has an omega-limit
point $z$ on the boundary, then its zero-coordinate set
$Z=\{i:z_i=0\}$ is a siphon.

\emph{Proof.}
The omega-limit set of a trajectory contained in a compact set is invariant
also in the following local backward sense: for sufficiently small
$s>0$, the solution through any omega-limit point can be followed back
to time $-s$ while staying in the nonnegative orthant. To see this,
choose $t_n\to\infty$ with $x(t_n)\to z$ and pass to a convergent
subsequence of $x(t_n-s)$. Its limit $z_{-s}$ is again an omega-limit
point and satisfies $\Phi_s(z_{-s})=z$ by continuity of the flow.
Local uniqueness identifies this with the backward solution through $z$.

Let $f$ denote the polynomial vector field. For $i\in Z$,
$f_i(z)\geq0$ by invariance of the orthant. If $f_i(z)>0$, then the
local backward solution has coordinate
$z_i(-s)=-s f_i(z)+o(s)<0$ for small positive $s$, contradicting the
preceding property. Thus $f_i(z)=0$ for every $i\in Z$.

Now consider a reaction whose reactant has no species in $Z$ but whose
product has $y'_i>0$ for some $i\in Z$. All the reactant's coordinates
at $z$ are positive, so its rate $k z^y$ is positive, including when
$y=0$. Its contribution to $f_i(z)$ is $kz^y y'_i>0$.
Every other nonzero contribution to that coordinate is nonnegative,
since reactions consuming $i$ vanish there. This contradicts
$f_i(z)=0$. No such reaction exists, which is exactly the siphon
condition.\hfill$\square$

This proof uses an omega-limit point, not an arbitrary boundary point.
At an arbitrary boundary point a positive inward derivative is perfectly
possible. The backward-invariance argument is what makes the stronger
zero-derivative conclusion legitimate here.

\paragraph{A complete structural sufficient condition.}
Say that a siphon $Z$ contains a conserved support if there is a nonzero
vector $c\geq0$ with
\[
 c\in S^\perp,
 \qquad \operatorname{supp}(c)\subseteq Z.
\]
Such vectors are often called nonnegative conservation vectors or
$P$-semiflows. They may have zero entries outside their support; the
strictly positive vector in Lemma 481.1 serves a different purpose.

\textbf{Proposition 481.3.} If every siphon contains a conserved support,
then every positive trajectory in a bounded compatibility class is
persistent. Weak reversibility is not needed for this sufficient criterion.

\emph{Proof.}
If persistence failed, some sequence of times tending to infinity would
have a coordinate tending to zero. Compactness gives a subsequence
converging to a boundary omega-limit point $z$. Lemma 481.2 makes its
zero set $Z$ a siphon. Choose the vector $c$ supplied by the hypothesis.
Since $x_0>0$ and $c\geq0$ is nonzero,
$c\cdot x_0>0$. Conservation and continuity imply
$c\cdot z=c\cdot x_0$, but the support condition gives $c\cdot z=0$,
a contradiction. Thus the omega-limit set lies in the positive orthant.
Its compactness gives a positive minimum of each coordinate on that
set; the trajectory's distance to its omega-limit set tends to zero,
so its coordinate liminf is also positive.\hfill$\square$

For the last distance assertion, a contrary sequence of times staying a
fixed distance away would have a convergent subsequence in the compact
class closure; its limit would, by definition, belong to the omega-limit
set. This supplies the claimed contradiction. The criterion is part of
the established siphon theory in [E2]; the proof above records its exact
use in the selected bounded-class setting.

Only finitely many subsets $Z$ need to be checked, and for a fixed $Z$
the conserved-support condition is a finite system of linear equalities
and inequalities, with a normalization such as $\sum c_i=1$.
It is enough to check minimal siphons: every other siphon contains a
minimal one, and its conserved support is still contained in the larger
set. This makes the criterion a practical certificate for particular
networks, but it is not automatically satisfied by every weakly
reversible network.

\paragraph{A critical siphon that still persists.}
Consider the reversible network
\[
 2A\ \mathop{\rightleftharpoons}^{k_1}_{k_2}\ A+B.
\]
It conserves $a+b=M>0$, and the singleton $Z=\{A\}$ is a siphon:
both reactant complexes already contain $A$. Yet every conservation
vector is a multiple of $(1,1)$, so no nonzero nonnegative one is
supported on $\{A\}$. This is a critical siphon, in the sense that
the conservation certificate fails.

The system nevertheless has a complete elementary persistence proof:
\[
 \dot a=-k_1a^2+k_2ab
        =a\bigl(k_2M-(k_1+k_2)a\bigr),\qquad b=M-a.
\]
For $0<a(0)<M$, let $a_*=k_2M/(k_1+k_2)$. The explicit solution is
\[
 a(t)=\frac{a_*}{1+(a_*/a(0)-1)e^{-k_2Mt}},
 \qquad
 b(t)\longrightarrow\frac{k_1M}{k_1+k_2}>0.
\]
The denominator is positive for all $t\geq0$, and $a(t)\to a_*>0$.
The example shows exactly why ``weak reversibility supplies a conserved
support for every siphon'' is a false intermediate claim. Critical
siphons require a dynamical repulsion argument rather than only a linear
invariant. The same two-species network also appears as a useful
mass-action benchmark in the supplied source [S1].

\paragraph{A full one-dimensional stoichiometric proof.}
\textbf{Proposition 481.4.} If the network is weakly reversible,
$\dim S\leq1$, and the positive compatibility class is bounded, then
every positive trajectory in that class is persistent, for all positive
rate constants.

\emph{Proof.}
If $S=0$, the trajectory is constant. Otherwise choose a nonzero vector
$v$ spanning $S$, and parametrize the class as
\[
 x(s)=x_0+s v,\qquad a<s<b.
\]
Boundedness makes both endpoints finite. At the lower endpoint define
$Z=\{i:x_i(a)=0\}$. It is nonempty, and $v_i>0$ for every $i\in Z$,
because $x_i(a+z)=v_i z>0$ for small $z>0$. Put
$V_Z=\sum_{i\in Z}v_i>0$.

For each reaction $r:y\to y'$ write $y'-y=c_r v$. Along the class,
\[
 \dot s=f(s)=\sum_r k_r x(s)^y c_r.
\]
Ignore self-reactions, which have $c_r=0$ and do not affect this
polynomial. For a complex $y$, put $m_y=\sum_{i\in Z}y_i$.
The exact linear vanishing of the coordinates in $Z$ gives
\begin{equation}\label{ra481:leading}
 x(a+z)^y=C_y z^{m_y}(1+O(z)),
 \quad
 C_y=\prod_{i\in Z}v_i^{y_i}
                   \prod_{i\notin Z}x_i(a)^{y_i}>0.
\end{equation}
The formula includes zero exponents in the usual way.

In any nontrivial linkage class, all differences of complexes are
parallel to $v$: connect the complexes by paths in the underlying
reaction graph and add the reaction differences. Moreover, distinct
complexes in that linkage class have distinct $m_y$, because
$m_{y'}-m_y=c V_Z$ if $y'-y=c v$. Hence there is a unique complex
of minimum $m_y$ in that linkage class. Weak reversibility makes the
linkage class strongly connected, so this minimum has at least one
outgoing non-self reaction. Every such reaction has $c_r>0$.
Every reaction with $c_r<0$ has a source with strictly larger $m_y$
than the minimum in its own linkage class.

Take the smallest of those minimum exponents over all nontrivial linkage
classes, call it $m_*$. All reactions with source exponent $m_*$ have
$c_r>0$, and at least one exists. All negative terms in $f(a+z)$ have
higher exponent. Since there are finitely many reactions,
\[
 f(a+z)=C_* z^{m_*}+O(z^{m_*+1}),\qquad C_*>0.
\]
Thus $f(s)>0$ near the lower endpoint. Apply the identical argument at
the upper endpoint with coordinate $b-s$ and direction $-v$; it gives
$f(s)<0$ near $b$.

Choose $\delta>0$ so that these signs hold in the endpoint intervals.
The closed interval
\[
 [\min\{s(0),a+\delta\},\ \max\{s(0),b-\delta\}]
 \ \subset\ (a,b)
\]
is forward invariant by the signs of the scalar vector field at its
endpoints and uniqueness of solutions. Every coordinate of $x(s)$ has
a positive minimum on that compact interval, proving persistence.
\hfill$\square$

This argument is a direct low-dimensional benchmark, within the already
known territory represented by [E3]. It does not claim that higher
dimensional classes reduce to one-dimensional ones. Its decisive feature
is that all vanishing coordinates are fixed positive multiples of the
same small parameter $s-a$.

\paragraph{Weak reversibility does not supply complex balance.}
Consider a second reversible network, with rates written on the arrows:
\[
 A\mathop{\rightleftharpoons}^{1}_{1}B,
 \qquad
 2A\mathop{\rightleftharpoons}^{1}_{4}2B.
\]
Every class $a+b=M$ is bounded and has stoichiometric dimension one, so
Proposition 481.4 applies. Complex balance at a positive point $(a,b)$
would require flow balance separately at each complex. The first
two-complex linkage class requires $a=b$; the second requires
$a^2=4b^2$, hence $a=2b$. These are incompatible for a positive point.
Thus this weakly reversible mass-action system is not complex-balanced
for its specified rates, even though it is persistent by the proof above.

The usual relative-entropy expression cannot simply be declared a
Lyapunov function with an arbitrary reference point. On the class
$a+b=1$, choose the reference $c=(1/2,1/2)$ and define
\[
 V_c(a,b)=a\log(2a)-a+\tfrac12
                  +b\log(2b)-b+\tfrac12.
\]
Here $\dot a=-a+b-2a^2+8b^2$. At $(a,b)=(3/5,2/5)$,
\[
 \dot a=\frac9{25},\qquad
 \dot V_c=\dot a\log(a/b)=\frac9{25}\log(3/2)>0.
\]
This calculation does not deny the possible existence of another
Lyapunov function. It shows that the complex-balanced entropy argument
cannot be imported without verifying its balance hypothesis. The recent
manuscript claim discussed at the start concerns a broader invariant-region
construction, not an assertion that all these rates are complex-balanced.

\paragraph{The first gap in extending the elementary proof.}
In higher-dimensional classes, several coordinates may vanish at
unrelated rates near a critical siphon face. For example, the two
monomials $x_1^2$ and $x_2$ have ratio $x_1^2/x_2$. Along
$x_1=x_2^3\to0$ the ratio tends to zero; along
$x_2=x_1^3\to0$ it tends to infinity. A ranking by total degree in
the vanishing species is therefore not a uniform ranking of reaction
rates. The one-parameter expansion \eqref{ra481:leading} is no longer
available along every approach to the face.

One would need a family of inward barriers that works across all such
relative-rate regimes, including transitions between them, or another
argument excluding critical siphons as omega-limit zero sets. Merely
having a directed cycle of complexes does not bound the instantaneous
ratios of its monomial rates. Nor does a conserved positive total mass
bound each summand away from zero. Proposition 481.3 handles faces with
a supported conservation law, and Proposition 481.4 handles one parameter;
neither supplies the required general barrier.

Dropping weak reversibility would make the target false even on a
bounded class: for $A\to B$ with rate $k a$, one has
$a(t)=a(0)e^{-kt}\to0$ and $a(t)+b(t)$ constant. This verifies that
finite-time positivity plus a compact class cannot replace the reaction
graph condition. Conversely, the reversible examples above show that
the graph condition can enforce persistence by nonlinear dynamics even
when the simplest conserved-support test fails.

\paragraph{Finite checks and outcome.}
The companion exact-arithmetic script enumerates all nonempty species
subsets in the displayed networks and tests the siphon condition directly
against the reaction list. It verifies that $2A\rightleftharpoons A+B$
has the critical siphon $\{A\}$, that the monomolecular directed cycle
$A\to B\to C\to A$ has only the full-species siphon, and that the
inconsistent-balance reversible example has the stated one-dimensional
vector field. Its rational calculations check the positive entropy
derivative and the logistic coefficients. No finite-time numerical
trajectory is used to certify an asymptotic positive liminf.

\textbf{Outcome of the independent argument: UNRESOLVED.} The work
proves the boundedness/conservation equivalence needed for the ODE setup,
the boundary-siphon lemma, the conserved-support persistence criterion,
and the full one-dimensional weakly reversible subcase. It also gives
explicit counterexamples to two intermediate assertions: weak reversibility
does not force every siphon to carry a conservation law, and it does not
force complex balance for arbitrary rate constants.

The remaining gap in this argument is a general boundary-exclusion
mechanism for critical siphons in higher-dimensional classes. Specific
next steps are to organize monomial dominance by relative logarithmic
rates, construct compatible inward barriers through the resulting
regimes, or independently audit the corresponding invariant-region
construction in the September 2026 manuscript [E1]. That manuscript
explicitly claims the bounded-trajectory conclusion covering this target;
this notebook neither verifies its entire proof nor claims to refute it.
The local attempt status must therefore not be read as a certified
current-literature verdict that the selected formulation remains open.

\subsection{Sources}
\begin{itemize}
\item[S1]Persistence and the Global Attractor Conjecture. Accessed 2026-09-01.
\url{https://reaction-networks.net/wiki/Persistence_and_the_Global_Attractor_Conjecture}.
\textit{Formulation source.}
\item[E1] G. Craciun, \emph{Toric Differential Inclusions and a Proof of
the Global Attractor Conjecture}, arXiv:1501.02860v3,
23 September 2026, Introduction, Theorem A / Theorem 3.1,
and the final remark in Section 8.
\url{https://arxiv.org/html/1501.02860v3}.
\textit{Very recent primary manuscript claiming bounded-trajectory
persistence for variable-$k$ weakly reversible systems; exact scope and
version checked 27 September 2026, full proof not independently audited.}
\item[E2] D. Angeli, P. De Leenheer, and E. D. Sontag,
\emph{A Petri Net Approach to Persistence Analysis in Chemical Reaction
Networks}, in Biology and Control Theory: Current Challenges,
Lecture Notes in Control and Information Sciences 357 (2007), 181--216,
Sections 4--5 on siphons and sufficient conservation conditions.
\url{https://sites.science.oregonstate.edu/~deleenhp/publication-list/toulouse.pdf}.
\textit{Primary author-hosted text; consulted 27 September 2026.}
\item[E3] C. Pantea, \emph{On the persistence and global stability of
mass-action systems}, arXiv:1103.0603,
abstract and the persistence theorem for two-dimensional stoichiometric
subspaces.
\url{https://arxiv.org/abs/1103.0603}.
\textit{Primary author record used only to identify established low-dimensional
scope; consulted 27 September 2026.}
\end{itemize}
\end{document}
